\chapter{内存、总线、I/O 与整机}

最小 CPU 能执行编码，但它还不是整机。整机需要更大的存储空间、部件间通信路径、设备控制方式、异步事件入口和复位启动路径。本章只站在硬件角度收束这些结构：数据怎样离开 CPU，怎样到达主存或设备，结果又怎样回来。

\section{一、主存扩展状态空间}

寄存器阵列很快，但容量有限。若所有状态都放在寄存器里，编号、连线、读写端口都会迅速膨胀。主存提供大容量状态空间，用地址选择位置，用数据线传递内容，用控制信号说明读还是写。

\begin{longtable}{p{0.2\linewidth}p{0.5\linewidth}}
\toprule
信号 & 作用 \\
\midrule
地址 & 指出访问哪个存储位置 \\
写数据 & 写入时提供数据 \\
读数据 & 读取时返回数据 \\
读写控制 & 指明本次访问方向 \\
ready/valid & 指明请求或响应是否有效 \\
\bottomrule
\end{longtable}

一次读访问可以这样理解：CPU 给出地址和读控制，主存经过内部译码、阵列访问和输出驱动后，把数据放到读数据线上，并用响应信号表示结果有效。一次写访问则是 CPU 同时给出地址、写数据和写控制，主存在合适边界把数据保存到对应位置。

主存和寄存器的差别不只是容量。寄存器通常直接接在 CPU 数据通路里，延迟短、端口少而固定；主存容量大，但访问路径更长，需要地址译码、阵列选择、总线传输和握手。

\section{二、快慢差距逼出层级存储}

CPU 内部一个周期可以很短，而主存访问可能慢得多。若每次取编码、读数据、写数据都等主存完成，ALU、寄存器阵列和控制器会长时间空转。层级存储的动机，就是把最近、最常用的一小部分内容放到更近、更快的位置。

cache 可以先理解成靠近 CPU 的小型硬件存储。它不是一句“先临时放一下”的比喻，而是一组硬件表项：

\begin{longtable}{p{0.2\linewidth}p{0.56\linewidth}}
\toprule
字段 & 作用 \\
\midrule
tag & 判断当前地址是否对应这一项 \\
data & 保存从主存取来的数据块 \\
valid & 表示这一项是否有效 \\
dirty & 表示这一项是否被改过、尚未写回下层 \\
\bottomrule
\end{longtable}

命中时，数据从 cache 返回，CPU 不必等主存。未命中时，硬件要向下层请求数据块，再填入 cache。若被替换的块是 dirty，还要先写回。这里先抓住硬件链路：cache 解决的是速度差距，但它引入了命中判断、替换、写回和一致性边界。

用一次读数据看这条链路。CPU 内部的数据通路算出地址 A，并发出读请求。地址 A 先被拆成几部分：低位选择块内的哪几个字节，中间位选择 cache 的某一组，高位作为 tag 和表项里保存的 tag 比较。若这一组中存在 valid=1 且 tag 相等的表项，就是命中。此时数据阵列读出整块数据，再由块内偏移选出需要的字节或字，送回 CPU 的读数据路径。对 CPU 来说，这次访问很快结束；对硬件来说，实际发生的是索引、tag 比较、有效位判断、数据选择和响应返回。

未命中时路径明显更长。cache 发现没有匹配表项后，不能凭空产生数据，只能向下层主存发出读块请求。若准备替换的旧表项 dirty=0，可以直接覆盖；若 dirty=1，说明旧块被改过但下层还没有最新内容，cache 必须先把旧块地址和旧数据写回主存，等写回完成后再读取新块。新块返回后，cache 更新 data、tag、valid，必要时清除 dirty，然后把 CPU 本来要的那一部分数据返回。这个过程解释了为什么同样是读一个地址，命中和未命中的硬件代价完全不同。

写访问还要再区分。若写命中，cache 可以直接修改对应字节，并把 dirty 置 1，表示这份内容已经比下层新；也可以同时写下层，保持下层同步。若写未命中，硬件策略会决定是先把整块取入 cache 再改，还是绕过 cache 直接写下层。不同策略会影响性能和复杂度，但核心问题一样：每一层都必须清楚“哪个副本是最新的、什么时候需要写回、CPU 何时可以认为写入完成”。

\begin{pdfdiagram}[x=1cm,y=1cm]
  \node[hw state,minimum width=1.45cm] (cpu) at (0.8,1.4) {CPU};
  \node[hw memory,minimum width=1.45cm] (l1) at (3.0,1.4) {cache};
  \node[hw memory,minimum width=1.45cm] (mem) at (5.25,1.4) {主存};
  \node[hw memory,minimum width=1.45cm] (storage) at (7.55,1.4) {外部存储};
  \draw[hw bus] (cpu.east) -- node[above,hw label] {最快/最小} (l1.west);
  \draw[hw bus] (l1.east) -- node[above,hw label] {较慢/较大} (mem.west);
  \draw[hw bus] (mem.east) -- node[above,hw label] {更慢/更大} (storage.west);
  \node[hw label,align=center] at (4.2,0.35) {层级存储用更近的小容量硬件吸收大量近期访问，未命中才向下层请求};
\end{pdfdiagram}
\diagramnote{层级存储不是单个部件，而是用速度、容量和距离形成的硬件链路。}

\section{三、总线和互连连接多个目标}

当 CPU、cache、主存、设备控制器都要通信时，需要互连结构。简单系统可以用共享总线，复杂系统会用分层互连或片上网络。无论形式如何，它至少要解决五个问题：

\begin{itemize}
\item 谁发起请求。
\item 地址指向哪个目标。
\item 数据往哪边走。
\item 多个发起者冲突时谁先走。
\item 响应怎样回到发起者。
\end{itemize}

共享总线的直觉是：多个部件共用一组地址线、数据线和控制线。由于同一时刻不能让多个源同时驱动总线，就需要仲裁器决定谁获得使用权。分层互连则把不同目标挂在不同分支上，通过译码和路由把请求送到正确目标。

互连不是“几根线”这么简单。它要保证电气上不冲突，时序上能握手，地址上能译码，响应上能返回正确发起者。整机复杂度很大一部分来自这里。

看一个最简单的互连访问。CPU 发出一个读请求，地址落在主存范围。互连先接收请求信号、地址和发起者编号，然后地址译码逻辑判断这个地址应该送到主存控制器，而不是设备控制器。若此时总线空闲，仲裁器授权 CPU，地址和控制信号被送到主存侧；若 DMA 正在占用总线，CPU 的请求可能要等待。主存返回数据时，互连还要把响应送回原来的发起者。若系统里只有 CPU 一个发起者，这件事看起来很简单；一旦 DMA、图形设备或其他主设备也能发请求，响应通道就必须带着足够的信息，保证数据不会回错对象。

再看地址落在设备范围的情况。CPU 发出的仍然是“地址、读写方向、数据、有效信号”这类硬件请求，但互连译码后不会把它送到主存，而是送到某个设备控制器。设备控制器看到的是自己的寄存器地址：可能是状态寄存器，可能是命令寄存器，也可能是数据缓冲寄存器。CPU 访问主存和访问设备，在数据通路入口处形式很像；真正的分岔发生在互连的地址译码处。

\begin{pdfdiagram}[x=1cm,y=1cm]
  \node[hw state,minimum width=1.45cm] (cpu) at (0.8,2.1) {CPU};
  \node[hw compute,minimum width=1.45cm] (dma) at (0.8,0.8) {DMA};
  \node[hw control,minimum width=1.65cm] (arb) at (3.2,1.45) {仲裁/互连};
  \node[hw memory,minimum width=1.45cm] (mem) at (5.65,2.35) {主存};
  \node[hw control,minimum width=1.55cm] (dev) at (5.65,1.15) {设备控制器};
  \node[hw control,minimum width=1.55cm] (intc) at (5.65,-0.05) {中断控制器};
  \draw[hw bus] (cpu.east) -- (arb.west);
  \draw[hw bus] (dma.east) -- (arb.west);
  \draw[hw bus] (arb.east) -- (mem.west);
  \draw[hw bus] (arb.east) -- (dev.west);
  \draw[hw bus] (dev.south) -- (intc.north);
  \draw[BookRed,line width=0.55pt,-{Latex[length=2mm,width=1.5mm]}] (intc.west) .. controls (3.1,-0.85) and (0.8,-0.35) .. node[below,hw label] {事件请求} (cpu.south);
  \node[hw label,align=center] at (3.45,-1.25) {总线/互连处理请求去向；中断控制器处理设备事件回到 CPU 的入口};
\end{pdfdiagram}
\diagramnote{整机互连同时要处理访问请求、设备寄存器、DMA 搬运和中断事件。}

\section{四、设备通过寄存器、DMA 和中断接入}

设备控制器通常把自己的控制寄存器映射到地址空间中。CPU 对某些地址读写，实际访问的是设备寄存器，这就是 MMIO。写一个控制寄存器可能启动设备动作，读一个状态寄存器可能得到设备是否完成、是否出错、缓冲区是否可用。

若设备要搬运大量数据，让 CPU 一字一字读写会浪费内部数据通路时间。DMA 控制器可以在获得总线使用权后，直接在设备和主存之间搬运数据。CPU 只需要配置源地址、目标地址、长度和方向，真正的数据搬运由 DMA 硬件完成。

设备完成工作或状态变化时，需要把事件送回 CPU。中断控制器负责收集多个设备的请求，进行屏蔽、优先级排序和投递。站在本书角度，中断先看成硬件事件入口：外部设备把一条请求线送进控制结构，CPU 在规定边界看到这个请求，并转入对应处理路径。

先看一次 MMIO 写。假设某个设备控制器有一个命令寄存器，地址是 D。CPU 执行一次写访问，地址线给出 D，写数据线上放入命令 bit，写控制有效。互连把 D 译码到这个设备控制器。设备控制器内部不是普通主存阵列，而是一组控制寄存器和状态机；它在握手边界接收写入值，把命令寄存器更新为新值，并可能让内部状态机从 idle 进入 busy。若命令要求设备开始工作，真正的机械、电气或外部协议动作会在设备控制器内部继续推进。CPU 写的是一个地址，硬件实际改变的是设备控制器的状态。

再看一次 MMIO 读。CPU 读取设备状态地址 S，互连把请求送到设备控制器。控制器把当前状态寄存器的值放到读数据线上，例如 busy 位、done 位、error 位、buffer\_ready 位。这个读操作不一定只是“看一眼”。有些设备寄存器会在读取后清除某些事件位，有些寄存器要求按特定顺序访问，否则会读到旧状态或破坏握手。因此设备寄存器虽然映射在地址空间里，语义却由设备控制器硬件定义。

DMA 的链路更能体现整机配合。若不用 DMA，CPU 要把一大块设备数据搬到主存，通常需要反复执行“读设备数据寄存器、写主存地址、地址递增、计数递减”这样的动作。每搬一个字，CPU 的寄存器、ALU、互连和设备寄存器都要参与一次。若用 DMA，CPU 先通过 MMIO 配置 DMA 控制器：源端是设备数据端口，目标是主存起始地址，长度是 N，方向是设备到主存，然后写启动位。之后 DMA 控制器自己成为总线发起者：它向设备侧取数据，向主存侧写数据，内部地址计数器递增，长度计数器递减，直到 N 变成 0。CPU 数据通路不再逐字搬运，但互连、主存控制器和设备控制器仍然在硬件层面真实传输每一个数据块。

DMA 也带来新的硬件边界。因为 CPU 和 DMA 都可能访问主存，互连必须仲裁；因为 CPU 可能从 cache 里看见旧数据，系统必须规定 DMA 写入的主存区域怎样和 cache 保持一致；因为设备可能比主存慢，DMA 控制器还要处理等待、重试或分批传输。这里不需要先引入操作系统，只要看到硬件事实：一旦有多个发起者同时读写共享存储，仲裁、可见性和完成信号就变成整机必须解决的问题。

最后看中断。设备完成 DMA 后，可以把 done 位写进自己的状态寄存器，同时拉起一条中断请求线。中断控制器接收来自多个设备的请求，先检查屏蔽位，再按优先级选择一个请求，把它整理成 CPU 能看到的事件信号。CPU 不会在任意组合逻辑中间突然改变动作，它通常在指令边界或规定的状态边界采样这个事件。采样到有效中断后，控制结构选择一条入口路径：保存必要状态，把 PC 改到事件入口地址，后续再由那里的编码读取设备状态并处理结果。

这条链路可以和轮询对比。轮询时，CPU 反复通过 MMIO 读取设备状态寄存器，直到 done 位变成 1；中断时，CPU 不必一直读状态寄存器，而是由设备完成后主动通过中断控制器发出事件。两者都建立在同一套硬件事实之上：设备状态保存在设备控制器里，CPU 通过互连读写这些寄存器，中断只是额外提供了一条“设备事件回到 CPU 控制入口”的硬件通道。

\section{五、复位启动给整机第一步}

通电或复位后，CPU、主存控制器、互连和设备控制器都必须进入规定初始状态。PC 通常被置到固定启动地址，那里连接 ROM、片上存储或其他启动介质。启动路径的第一件事，是让硬件从可预测位置开始取编码。

整机启动不仅是 CPU 复位。时钟源要稳定，复位树要按顺序释放，主存控制器要完成训练或初始化，互连要进入可转发请求的状态，必要设备要处于安全默认值。若这些硬件条件没有满足，即使 CPU 本身能取指，也可能在第一次访问主存或设备时卡住。

到这里，本书的硬件链路闭合：电和器件形成可控开关，开关形成门，门形成组合逻辑，反馈和时钟形成状态，寄存器和 ALU 形成数据通路，控制器组织动作，最小 CPU 执行编码，主存、cache、互连、设备和启动路径把它扩展成整机。

\begin{classicnote}
整机硬件的标注对象是事务。一次访问不是“读内存”三个字，而是地址、方向、数据、发起者、目标、握手、响应和错误状态组成的一次硬件事务。

\begin{itemize}
\item 纸上实验：分别走读一次 cache 命中读、cache 未命中读、MMIO 写设备命令、DMA 完成并发中断，标出每一步的发起者和响应者。
\item 机器证据：主存、cache、设备寄存器都可以被地址访问，但语义不同；地址译码决定请求落到哪类目标。
\item 并发证据：DMA 和 CPU 都能成为总线发起者，必须有仲裁；cache 和 DMA 共享主存时，必须有一致性或同步规则。
\item 检查题：CPU 写设备命令寄存器后，能否立刻假设设备动作完成？应该看 MMIO 写完成、设备 done 位、中断，还是 DMA completion？
\end{itemize}

经典系统教材会把整机看成多层抽象：指令看到地址，硬件看到事务，设备看到寄存器和状态机，软件最终看到完成事件。能把这几层对应起来，才算从最小 CPU 走到了真正机器。
\end{classicnote}
